\documentclass[11pt,a4paper]{article}
\usepackage[margin=1in]{geometry}
\usepackage{amsmath,amssymb,amsthm,bm}
\usepackage{booktabs,array,multirow}
\usepackage{hyperref}
\usepackage{natbib}
\hypersetup{colorlinks=true,linkcolor=blue,citecolor=blue}
\usepackage{xcolor}
\usepackage{listings}
\lstset{
  language=Matlab,
  basicstyle=\small\ttfamily,
  keywordstyle=\color{blue},
  commentstyle=\color{gray},
  stringstyle=\color{orange},
  breaklines=true,
  frame=single,
  numbers=left,
  numberstyle=\tiny\color{gray},
  xleftmargin=1.5em,
  framexleftmargin=1.5em
}

% ── shortcuts ──────────────────────────────────────────────────────────────
\newcommand{\E}{\mathbb{E}}
\newcommand{\OmH}{\bm\Omega^{H}}
\newcommand{\OmF}{\bm\Omega^{F}}
\newcommand{\lD}{\bm\lambda_D}
\newcommand{\wh}[1]{\widehat{#1}}
\providecommand{\tl}[1]{\widetilde{#1}}
\newcommand{\DHmat}{\bm\Delta}
\newcommand{\DC}{\pi^{DC}}

\newtheorem{proposition}{Proposition}
\newtheorem{lemma}{Lemma}

\title{\textbf{Open-Economy Production Network New Keynesian Model}\\[4pt]
\large Model Equations, Calibration, and Dynare Implementation}
\author{}
\date{\today}

\begin{document}
\maketitle
\tableofcontents
\bigskip

% ════════════════════════════════════════════════════════════════════════════
\section{Model Overview}
% ════════════════════════════════════════════════════════════════════════════

This model extends \citet{Rubbo2024} to a small open economy.  The closed-economy
model shows that, with multiple sectors and intermediate-input linkages, productivity
shocks generate an inflation–output tradeoff for every standard price index except
a unique \emph{divine coincidence} (DC) index, which weights sectors by
Domar weights scaled by price stickiness.  We add three open-economy
ingredients: (i)~imported intermediate inputs (with home/foreign Armington
nesting at each production node), (ii)~imported final consumption goods, and
(iii)~a single internationally traded bond with a debt-elastic premium.

The key result is that the exchange rate enters sectoral marginal costs exactly
like a productivity shock scaled by the import-content matrix $\OmF$.
By embedding import nodes as fully flexible \emph{upstream sectors}
($\wh\delta_{M}=1$) in an augmented network, the DC index still exists
and is immune to both domestic TFP shocks and FX/import-price shocks.
A float targeting the DC index implements zero output gap and fully
insulates the economy from external shocks via exchange-rate adjustment.
A hard peg forces the full adjustment onto domestic inflation and output.

The model is calibrated to a \textbf{three-sector oil-exporting economy}:
an upstream Resource sector (import-intensive inputs, most output exported),
a Manufacturing sector (uses resources as intermediates), and a Services sector
(sells mainly to domestic consumers).  Three shocks are studied---an import
price shock ($\varepsilon^{pF}$), a domestic TFP shock ($\varepsilon^{a}$),
and a foreign demand shock ($\varepsilon^{D}$)---under three policy regimes:
free float, hard peg, and managed float.

% ════════════════════════════════════════════════════════════════════════════
\section{Environment}
% ════════════════════════════════════════════════════════════════════════════

\subsection{Indexing and notation}

There are $N=3$ domestic sectors indexed $i,j\in\{1,2,3\}$
(Resource, Manufacturing, Services).
Lower-case hats denote log-deviations from the non-stochastic steady state;
tildes denote gaps (sticky-price minus flex-price).
$E_t[\cdot]$ is the rational expectation conditional on information at $t$.

\paragraph{Network matrices.}
\begin{align*}
\Omega^H_{ij}&=\text{domestic cost share of input }j\text{ in sector }i\text{'s costs},\\
\Omega^F_{ij}&=\text{imported cost share of input }j\text{ in sector }i\text{'s costs},\\
\alpha_i&=1-\sum_j(\Omega^H_{ij}+\Omega^F_{ij})=\text{labour share in sector }i.
\end{align*}
The domestic Leontief inverse is $\mathbf{L}^H\equiv(\mathbf{I}-\OmH)^{-1}$.
The \textbf{domestic supplier centrality} (open-economy Domar weight) is
\begin{equation}
\lambda_{D,i}\;\equiv\;\bigl[\bm\beta^{H\top}\mathbf{L}^H\bigr]_i,
\label{eq:lambdaD}
\end{equation}
where $\bm\beta^H=(\beta^H_1,\beta^H_2,\beta^H_3)^{\top}$ are domestic final
consumption shares.

\paragraph{Calvo adjustment.}
$\delta_i$ is the probability that a firm in sector $i$ resets its price each period.
Following \citet{Rubbo2024}, define the adjusted parameter
\begin{equation}
\wh\delta_i\;\equiv\;\frac{\delta_i\bigl(1-\beta(1-\delta_i)\bigr)}{1-\beta\delta_i(1-\delta_i)}.
\label{eq:deltahat}
\end{equation}

\subsection{Households}

The representative household maximises
\begin{equation}
\E_0\sum_{t=0}^{\infty}\beta^t\!\left[\frac{C_t^{1-\gamma}}{1-\gamma}-\frac{L_t^{1+\varphi}}{1+\varphi}\right].
\end{equation}
Aggregate consumption is a CES composite of a domestic bundle $C^H_t$ and an
imported bundle $C^F_t$:
\begin{equation}
C_t=\left[\omega^{1/\eta}(C^H_t)^{\frac{\eta-1}{\eta}}+(1-\omega)^{1/\eta}(C^F_t)^{\frac{\eta-1}{\eta}}\right]^{\frac{\eta}{\eta-1}},
\end{equation}
where $\omega$ is home bias in final consumption and $\eta$ is the
home–foreign elasticity.  The domestic bundle aggregates across sectors
with shares $\bm\beta^H$; the foreign bundle aggregates across $M$ imported
goods with shares $\bm\mu$ and elasticity $\sigma$.

The budget constraint is
\begin{equation}
P^C_tC_t+\frac{B_{t+1}}{1+i_t}+S_t\frac{B^*_{t+1}}{1+i^{*}_t}
\;\leq\; W_tL_t+\Pi_t-T_t+B_t+S_tB^*_t,
\end{equation}
where $S_t$ is the nominal exchange rate (home currency per unit of foreign),
$B_t$ is a domestic bond, $B^*_t$ is a foreign-currency bond, and
$\Pi_t$ are firm profits rebated lump sum.

The effective world rate includes a \textbf{debt-elastic premium}:
\begin{equation}
i^{*}_t\;=\;i^*-\psi\!\left(\frac{S_tB^*_t}{P^C_t}-\bar b^*\right),
\quad\psi>0,
\label{eq:premium}
\end{equation}
which ensures stationarity of net foreign assets (requires $\psi>1/\beta-1$).

\subsection{Firms}

Each domestic sector $i$ has a unit mass of monopolistically competitive firms
with CRS production:
\begin{equation}
Y_{ift}=A_{it}F_i\!\bigl(L_{ift},\{X_{ijft}\}_{j=1}^N\bigr).
\end{equation}
The intermediate input bundle $X_{ijft}$ is an Armington composite of domestic
and imported varieties of good $j$:
\begin{equation}
X_{ijft}=\left[v_{ij}^{1/\xi}(X^H_{ijft})^{\frac{\xi-1}{\xi}}+(1-v_{ij})^{1/\xi}(X^F_{ijft})^{\frac{\xi-1}{\xi}}\right]^{\frac{\xi}{\xi-1}}.
\end{equation}
Cost minimisation gives nominal marginal cost
\begin{equation}
MC_{it}=\frac{W_t^{\alpha_i}\prod_j\tilde P_{ijt}^{\omega_{ij}}}{A_{it}},
\label{eq:mc_structural}
\end{equation}
where $\tilde P_{ijt}$ is the Armington price of the input bundle (weighted
average of domestic price $P_{jt}$ and foreign price $S_tP^*_{F,jt}$).
Firms face Calvo pricing: fraction $\delta_i$ reset optimally each period.

\subsection{Policy and shocks}

The central bank sets the nominal interest rate $i_t$ under one of three regimes
(Section~\ref{sec:policy}).  Exogenous shocks are AR(1) processes.

% ════════════════════════════════════════════════════════════════════════════
\section{Log-Linearised System}
% ════════════════════════════════════════════════════════════════════════════

All variables are log-deviations from the zero-inflation steady state.
The full system for the three-sector model is listed below; these are
the equations entered directly into Dynare.

\subsection{Sectoral Phillips curves}

For each sector $i\in\{1,2,3\}$, the Calvo pricing condition yields:
\begin{equation}
\boxed{
\pi_{it}=\beta(1-\wh\delta_i)\,\E_t\pi_{i,t+1}
+\wh\delta_i\,\widetilde{mc}_{it}
}
\label{eq:PC_generic}
\end{equation}
where the \textbf{real marginal cost gap} is
\begin{equation}
\widetilde{mc}_{it}
=\alpha_i\wh w_t
+\sum_{j}\omega^H_{ij}\,p_{jt}
+\omega^F_i(\wh s_t+\wh p^{F*}_t)
-\wh a_{it}
-p_{i,t-1}
\label{eq:mctilde}
\end{equation}
with $\omega^F_i\equiv\sum_j\Omega^F_{ij}$ (sector $i$'s total import share
in costs) and $p_{jt}=p_{j,t-1}+\pi_{jt}$.

\paragraph{Sector 1 — Resource:} no domestic inputs ($\Omega^H_{1j}=0$),
30\% imported machinery.
\begin{equation}
\pi_{1t}=\beta(1-\wh\delta_1)\E_t\pi_{1,t+1}
+\wh\delta_1\!\left[\alpha_1\wh w_t+\omega^F_1(\wh s_t+\wh p^{F*}_t)-\wh a_{1t}-p_{1,t-1}\right]
\label{eq:PC1}
\end{equation}

\paragraph{Sector 2 — Manufacturing:} 20\% from domestic Resource, 10\% imported.
\begin{equation}
\pi_{2t}=\beta(1-\wh\delta_2)\E_t\pi_{2,t+1}
+\wh\delta_2\!\left[\alpha_2\wh w_t+\omega^H_{21}p_{1t}+\omega^F_2(\wh s_t+\wh p^{F*}_t)-\wh a_{2t}-p_{2,t-1}\right]
\label{eq:PC2}
\end{equation}

\paragraph{Sector 3 — Services:} 25\% from domestic Manufacturing, 5\% imported.
\begin{equation}
\pi_{3t}=\beta(1-\wh\delta_3)\E_t\pi_{3,t+1}
+\wh\delta_3\!\left[\alpha_3\wh w_t+\omega^H_{32}p_{2t}+\omega^F_3(\wh s_t+\wh p^{F*}_t)-\wh a_{3t}-p_{3,t-1}\right]
\label{eq:PC3}
\end{equation}

\subsection{Price level laws of motion}

\begin{equation}
p_{it}=p_{i,t-1}+\pi_{it},\qquad i=1,2,3.
\label{eq:pLOM}
\end{equation}

\subsection{Real wage (labour supply)}

Log-linearising the intratemporal optimality condition and using Lemma 2
(output gap $=$ employment gap):
\begin{equation}
\wh w_t=\wh p^C_t+(\gamma+\varphi)\,\tl y_t,
\label{eq:labsupply}
\end{equation}
where the \textbf{consumer price level} in log-deviations is
\begin{equation}
\wh p^C_t=\sum_i\beta^H_i\,p_{it}+\beta^{F,\text{tot}}\!(\wh s_t+\wh p^{F*}_t).
\label{eq:CPI_level}
\end{equation}

\subsection{CPI inflation}

\begin{equation}
\pi^C_t=\sum_i\beta^H_i\,\pi_{it}
+\beta^{F,\text{tot}}\!\bigl[(\wh s_t+\wh p^{F*}_t)-(\wh s_{t-1}+\wh p^{F*}_{t-1})\bigr].
\label{eq:CPIinfl}
\end{equation}

\subsection{Divine coincidence (DC) index}

\begin{equation}
\DC_t=\sum_i w_i\,\pi_{it},
\qquad
w_i=\frac{\lambda_{D,i}(1-\wh\delta_i)/\wh\delta_i}{\sum_k\lambda_{D,k}(1-\wh\delta_k)/\wh\delta_k}.
\label{eq:DCindex}
\end{equation}
By Proposition~\ref{prop:DC}, $\DC_t=\rho\E_t\DC_{t+1}+\kappa^{DC}\tl y_t$
with no residual term.

\subsection{IS curve}

Log-linearising the Euler equation and imposing market clearing:
\begin{equation}
\tl y_t=\E_t\tl y_{t+1}-\frac{1}{\gamma}\!\left(\wh\imath_t-\E_t\pi^C_{t+1}-r^{nat}_{t+1}\right)+\nu_D\wh D^*_t,
\label{eq:IS}
\end{equation}
where $r^{nat}_{t+1}$ is the natural (flex-price) real interest rate and
$\nu_D>0$ captures the direct effect of foreign demand on domestic output.

\subsection{Uncovered interest parity (UIP)}

\begin{equation}
\wh\imath_t=\wh\imath^*-\psi\,\wh b^*_t+\E_t\wh s_{t+1}-\wh s_t.
\label{eq:UIP}
\end{equation}

\subsection{Net foreign assets (NFA)}

The balance-of-payments identity in log-linear form:
\begin{equation}
\wh b^*_{t+1}=\underbrace{\left(\frac{1}{\beta}-\psi\right)}_{\displaystyle<1}\wh b^*_t
+\kappa_{EX}\!\left(-\theta^*\wh s_t+\wh D^*_t\right)
-\kappa_{IM}\!\left(\wh p^{F*}_t+\tl y_t\right),
\label{eq:NFA}
\end{equation}
where $\kappa_{EX}$ and $\kappa_{IM}$ are steady-state ratios of exports and
imports to GDP.  The eigenvalue $(1/\beta-\psi)<1$ ensures stationarity.

\subsection{Exchange rate pass-through to sectoral costs}

Defining \textbf{import centrality}
\begin{equation}
\mathcal{M}_i\;\equiv\;\bigl[(\mathbf{I}-\OmH)^{-1}\OmF\mathbf{1}\bigr]_i,
\label{eq:importcent}
\end{equation}
the total exposure of sector $i$'s inflation to an FX/import-price shock
under a peg is proportional to $\mathcal{M}_i$.  Sectors with higher import
centrality transmit more FX shocks to domestic prices.

\subsection{Policy rules}
\label{sec:policy}

Three alternative regimes are implemented by replacing one block in Dynare:

\paragraph{Regime A — Free float (benchmark):}
\begin{equation}
\wh\imath_t=\phi_\pi\DC_t+\phi_y\tl y_t.
\label{eq:Taylor_float}
\end{equation}
$\wh s_t$ is endogenous (determined by UIP).

\paragraph{Regime B — Hard peg:}
\begin{equation}
\wh s_t=0\quad\forall\,t.
\label{eq:peg}
\end{equation}
$\wh\imath_t$ is residually determined by UIP \eqref{eq:UIP}.

\paragraph{Regime C — Managed float:}
\begin{equation}
\wh\imath_t=\phi_\pi\DC_t+\phi_y\tl y_t+\phi_s\wh s_t.
\label{eq:Taylor_managed}
\end{equation}
$\phi_s>0$ implies partial exchange-rate stabilisation; $\phi_s\to\infty$
approaches a peg.

\subsection{Exogenous shock processes}

\begin{align}
\wh a_{it}&=\rho_a\,\wh a_{i,t-1}+\varepsilon^a_{it},&i&=1,2,3,\label{eq:TFP}\\
\wh p^{F*}_t&=\rho_{pF}\,\wh p^{F*}_{t-1}+\varepsilon^{pF}_t,\label{eq:impPrice}\\
\wh D^*_t&=\rho_D\,\wh D^*_{t-1}+\varepsilon^D_t.\label{eq:demand}
\end{align}

\subsection{Variable summary}

\begin{center}
\begin{tabular}{lll}
\toprule
\textbf{Variable} & \textbf{Type} & \textbf{Description}\\
\midrule
$\pi_{1t},\pi_{2t},\pi_{3t}$ & Jump & Sectoral inflation (Resource, Manuf., Services)\\
$\tl y_t$ & Jump & Output gap\\
$\wh s_t$ & Jump (float) / Exog.=0 (peg) & Nominal exchange rate (dep.$>$0)\\
$\wh\imath_t$ & Jump (float) / Residual (peg) & Domestic interest rate\\
$p_{1t},p_{2t},p_{3t}$ & State (endogenous) & Sectoral price levels\\
$\wh b^*_t$ & State (endogenous) & Net foreign assets\\
$\wh w_t$ & Static & Nominal wage\\
$\pi^C_t,\wh p^C_t,\DC_t$ & Static & CPI inflation, CPI level, DC index\\
$\wh a_{1t},\wh a_{2t},\wh a_{3t}$ & State (exogenous) & Sectoral TFP\\
$\wh p^{F*}_t$ & State (exogenous) & Import price shock\\
$\wh D^*_t$ & State (exogenous) & Foreign demand shock\\
\bottomrule
\end{tabular}
\end{center}

% ════════════════════════════════════════════════════════════════════════════
\section{Key Analytical Results}
% ════════════════════════════════════════════════════════════════════════════

\begin{lemma}[Output gap and markups]
$(\gamma+\varphi)\tl y_t=-\sum_i\lambda_{D,i}\mu_{it}$,
where $\mu_{it}=p_{it}-mc_{it}$ is sector $i$'s log markup.
\end{lemma}

\begin{lemma}[Markup–inflation link]
$\mu_{it}=-\dfrac{1-\wh\delta_i}{\wh\delta_i}(\pi_{it}-\beta\E_t\pi_{i,t+1}).$
\end{lemma}

\begin{proposition}[Open-economy divine coincidence]
\label{prop:DC}
The DC index \eqref{eq:DCindex} satisfies
\begin{equation}
\DC_t=\beta\,\E_t\DC_{t+1}+\kappa^{DC}\tl y_t,
\quad
\kappa^{DC}=\frac{\gamma+\varphi}{\displaystyle\sum_k\lambda_{D,k}(1-\wh\delta_k)/\wh\delta_k},
\end{equation}
with \textbf{no residual}: neither domestic TFP shocks $\wh{\mathbf{a}}_t$
nor the exchange rate $\wh s_t$ nor import prices $\wh p^{F*}_t$ appear.

\emph{Proof sketch:} Embed each imported input as a fully-flexible upstream
node ($\wh\delta_M=1$) in an augmented network. Fully flexible nodes receive
zero weight in the DC index ($w_M=0$) because $(1-\wh\delta_M)/\wh\delta_M=0$.
Rubbo's Proposition~1 applies verbatim to the augmented network. $\square$
\end{proposition}

\paragraph{Corollary (FX regime comparison).}
Under the float \eqref{eq:Taylor_float} with $\phi_\pi>1$, the exchange rate
absorbs import-price shocks: $\wh s_t\approx-\wh p^{F*}_t$, so all sectoral
inflations converge to zero and $\tl y_t\to 0$.  Under the peg \eqref{eq:peg},
import-price shocks pass through into domestic inflation with magnitude
proportional to the import centrality $\mathcal{M}_i$ of each sector.

% ════════════════════════════════════════════════════════════════════════════
\section{Calibration}
% ════════════════════════════════════════════════════════════════════════════

\begin{table}[h!]
\centering
\caption{Calibration --- 3-sector oil-exporter (quarterly)}
\label{tab:calib}
\begin{tabular}{llll}
\toprule
\textbf{Parameter} & \textbf{Value} & \textbf{Description} & \textbf{Source}\\
\midrule
\multicolumn{4}{l}{\textit{Preferences}}\\
$\beta$ & 0.99 & Discount factor & Standard\\
$\gamma$ & 1.00 & Risk aversion / inverse EIS & Standard\\
$\varphi$ & 2.00 & Inverse Frisch elasticity & \citet{Smets2003}\\
$\eta$ & 1.50 & Home–foreign substitution & \citet{Gali2005}\\
$\omega$ & 0.70 & Home bias in consumption & Data average\\
\midrule
\multicolumn{4}{l}{\textit{Domestic input-output matrix $\OmH$}}\\
$\omega^H_{21}$ & 0.20 & Manuf.\ buys 20\% from Resource & BEA I-O\\
$\omega^H_{32}$ & 0.25 & Services buys 25\% from Manuf. & BEA I-O\\
\midrule
\multicolumn{4}{l}{\textit{Imported input matrix $\OmF$}}\\
$\omega^F_1$ & 0.30 & Resource: 30\% imported machinery & SOE data\\
$\omega^F_2$ & 0.10 & Manuf.: 10\% imported components & SOE data\\
$\omega^F_3$ & 0.05 & Services: 5\% indirect imports & SOE data\\
\midrule
\multicolumn{4}{l}{\textit{Labour shares $\alpha_i=1-\sum_j(\omega^H_{ij}+\omega^F_{ij})$}}\\
$\alpha_1$ & 0.70 & Resource & Derived\\
$\alpha_2$ & 0.70 & Manufacturing & Derived\\
$\alpha_3$ & 0.70 & Services & Derived\\
\midrule
\multicolumn{4}{l}{\textit{Final consumption shares}}\\
$\beta^H_1,\beta^H_2,\beta^H_3$ & 0.05,\ 0.15,\ 0.80 & Domestic sectors & Expenditure survey\\
$\beta^{F,\text{tot}}$ & 0.10 & Total import share in consumption & Trade data\\
\midrule
\multicolumn{4}{l}{\textit{Pricing (Calvo)}}\\
$\delta_1$ & 0.75 & Resource: flexible (commodity prices) & \citet{Pasten2017}\\
$\delta_2$ & 0.50 & Manufacturing: intermediate & \citet{Pasten2017}\\
$\delta_3$ & 0.25 & Services: sticky & \citet{Pasten2017}\\
$\varepsilon_i$ & 8.00 & Within-sector substitution elasticity & Standard NK\\
\midrule
\multicolumn{4}{l}{\textit{Open economy}}\\
$\psi$ & 0.020 & NFA debt-elastic premium (ensures stationarity) & \citet{SGU2003}\\
$\theta^*$ & 2.00 & Export price elasticity & \citet{Gali2005}\\
$\kappa_{EX}$ & 0.85 & Exports/GDP ratio & Oil exporter data\\
$\kappa_{IM}$ & 0.45 & Imports/GDP ratio & Oil exporter data\\
$\nu_D$ & 0.30 & Foreign demand sensitivity & Calibrated\\
\midrule
\multicolumn{4}{l}{\textit{Monetary policy}}\\
$\phi_\pi$ & 1.50 & Taylor rule: inflation & \citet{Taylor1993}\\
$\phi_y$ & 0.50 & Taylor rule: output gap & \citet{Taylor1993}\\
$\phi_s$ & 0.30 & Managed float: FX weight & Calibrated\\
\midrule
\multicolumn{4}{l}{\textit{Shock persistence}}\\
$\rho_a$ & 0.90 & TFP & Standard\\
$\rho_{pF}$ & 0.85 & Import price & Commodity price data\\
$\rho_D$ & 0.80 & Foreign demand & Standard\\
\midrule
\multicolumn{4}{l}{\textit{Derived network objects}}\\
$\lambda_{D,1},\lambda_{D,2},\lambda_{D,3}$ & 0.12,\ 0.35,\ 0.80 & Domestic supplier centrality & Eq.~\eqref{eq:lambdaD}\\
$w^{DC}_1,w^{DC}_2,w^{DC}_3$ & 0.005,\ 0.069,\ 0.926 & DC index weights & Eq.~\eqref{eq:DCindex}\\
$\kappa^{DC}$ & 0.298 & DC Phillips curve slope & Eq.~\eqref{prop:DC}\\
$\mathcal{M}_1,\mathcal{M}_2,\mathcal{M}_3$ & 0.30,\ 0.16,\ 0.09 & Import centrality & Eq.~\eqref{eq:importcent}\\
\bottomrule
\end{tabular}
\end{table}

The adjusted Calvo parameters \eqref{eq:deltahat} are
$(\wh\delta_1,\wh\delta_2,\wh\delta_3)=(0.693,\,0.336,\,0.079)$.

% ════════════════════════════════════════════════════════════════════════════
\section{Dynare Implementation}
% ════════════════════════════════════════════════════════════════════════════

\texttt{open\_economy\_network.mod} implements the \emph{original nonlinear}
structural model -- Cobb-Douglas cost minimization, the exact recursive
Calvo pricing problem, the CRRA/Frisch household FOCs, CES consumption
aggregation, and UIP with a debt-elastic premium -- not the
log-linearized system of Section~3. Dynare computes the non-stochastic
steady state numerically (\texttt{initval}/\texttt{steady;}) and then
perturbs around it (\texttt{stoch\_simul(order=2,pruning,...)}); the
log-linear system above is kept only as an independent theory benchmark
to sanity-check Dynare's own order-1 decision rules against. Two
modelling choices worth flagging:
\begin{itemize}
\item \textbf{Output gap.} Rather than solving a parallel flex-price
  shadow economy, $\tl y_t$ uses Lemma~1 in closed form: $\tl y_t =
  \sum_i\lambda_{D,i}\bigl[\log(Y_{it}/\bar Y_i)-\log A_{it}\bigr]$,
  an exact implication of the model rather than an approximation.
\item \textbf{$\wh\delta_i$ is not a primitive here.} The raw Calvo
  reset probabilities $\delta_i=(0.75,0.50,0.25)$ enter the nonlinear
  recursion directly; $\wh\delta_i$ (eq.~\ref{eq:deltahat}) is computed
  only to build the theory-implied DC-index weights fed into the
  Taylor rule, since it is itself a log-linearization artifact of the
  exact reset-price recursion (proofs.tex, Calvo section).
\end{itemize}

Regimes are switched via the \texttt{@\#define REGIME} macro at the top
of the file (\texttt{"float"}, \texttt{"peg"}, \texttt{"managed"}).
\texttt{code/run\_all\_regimes.m} runs all three automatically (by
generating one temporary \texttt{.mod} file per regime and calling
Dynare on each) and exports IRFs, the derived network objects
($\lambda_D$, import centrality, $\wh\delta$, DC weights) and order-2
variances to \texttt{results/*.csv}. \texttt{code/analysis.py} then
loads those CSVs to build the IRF, network/centrality, DC-insulation and
welfare-decomposition figures referenced in
\texttt{OpenEconomy\_Networks\_FX.pptx}; see that file's docstring for
the full pipeline. Superseded illustrative code that solved a hand-rolled
approximate linear system (not a genuine rational-expectations solution)
previously lived in \texttt{code/model.py}.

For reference, the file skeleton (see the repository for the full,
current version -- it changes independently of this write-up):

\begin{lstlisting}[caption={open\_economy\_network.mod --- structure only, see repo for current content},label={lst:dynare}]
// var Y1 Y2 Y3 L1 L2 L3 P1 P2 P3 PI1 PI2 PI3 MC1 MC2 MC3 PSTAR1-3
//     X1_i X2_i (Calvo recursion) C CH CF PH PC PIC Ltot W
//     S PF PFH BSTAR I EX IM DSTAR A1 A2 A3 piDC y_gap GDP;
// varexo eps_a1 eps_a2 eps_a3 eps_pF eps_D;
//
// model;
//   MC_i          = Cobb-Douglas nominal marginal cost
//   X1_i, X2_i    = exact Calvo recursive aggregators (real MC/P, not MC)
//   PSTAR_i       = optimal reset price = (eps/(eps-1))*X1_i/X2_i
//   1 = (1-delta_i)*PI_i^(eps-1) + delta_i*PSTAR_i^(1-eps)   [Calvo law]
//   L_i           = Cobb-Douglas labour demand
//   Y_i           = domestic C + intermediate demand + exports (mkt clearing)
//   PH, PC, CH, CF = Cobb-Douglas / CES consumption aggregation
//   Euler, labour supply (nonlinear CRRA/Frisch FOCs)
//   UIP with debt-elastic premium; NFA current-account identity
//   EX, IM        = export demand / total import demand
//   policy rule   = float | peg | managed (regime-dependent)
//   piDC, y_gap   = DC index (Prop. 1 weights); output gap (Lemma 1, closed form)
//   AR(1) shock processes in logs
// end;
\end{lstlisting}

\emph{(Retained below for historical reference: the very first
hand-linearized draft of this file, before it was rewritten as the
nonlinear structural model described above.)}

\begin{lstlisting}[caption={Superseded: original hand-linearized draft},label={lst:dynare-old}]
// =========================================================
// Open Economy Production Network NK Model
// Rubbo (2024) extended to Small Open Economy
// 3 sectors: Resource (1), Manufacturing (2), Services (3)
//
// REGIMES (uncomment one):
// @#define REGIME = "float"    // Free float (default)
// @#define REGIME = "peg"      // Hard peg
// @#define REGIME = "managed"  // Managed float
@#define REGIME = "float"
// =========================================================

// ---------------------------------------------------------
// VARIABLES
// ---------------------------------------------------------
var
    // Sectoral inflation
    pi1 pi2 pi3

    // Aggregate inflation measures
    piDC piCPI

    // Output gap and interest rate
    y_gap i_hat

    // Exchange rate (log deviation, depreciation = positive)
    s_hat

    // Sectoral price levels (log deviation)
    p1 p2 p3

    // Consumer price level (log deviation)
    pC_hat

    // Real wage (log deviation)
    w_hat

    // Net foreign assets (log deviation)
    bstar

    // TFP shocks (sectoral)
    a1 a2 a3

    // Import price shock
    pF

    // Foreign demand shock
    D_star
    ;

// ---------------------------------------------------------
// EXOGENOUS SHOCKS
// ---------------------------------------------------------
varexo eps_a1 eps_a2 eps_a3 eps_pF eps_D;

// ---------------------------------------------------------
// PARAMETERS
// ---------------------------------------------------------
parameters
    // Preferences
    BETA GAMMA VARPHI

    // Adjusted Calvo parameters (Rubbo eq. after (14))
    DHAT1 DHAT2 DHAT3

    // Labour shares
    ALPHA1 ALPHA2 ALPHA3

    // Domestic input-output shares
    OH21      // Manuf. buys from Resource (domestic)
    OH32      // Services buys from Manuf. (domestic)

    // Import shares (total per sector)
    OF1 OF2 OF3

    // Final consumption shares (domestic sectors)
    BH1 BH2 BH3

    // Total import share in final consumption
    BF_TOT

    // DC index weights
    WDC1 WDC2 WDC3

    // Open economy
    PSI       // NFA premium
    THETA_S   // export price elasticity
    KAP_EX    // exports/GDP
    KAP_IM    // imports/GDP
    NU_D      // foreign demand sensitivity

    // Policy
    PHI_PI PHI_Y PHI_S

    // Shock persistence
    RHO_A RHO_PF RHO_D
    ;

// ---------------------------------------------------------
// PARAMETER VALUES
// ---------------------------------------------------------

// Preferences
BETA    = 0.99;
GAMMA   = 1.00;
VARPHI  = 2.00;

// Adjusted Calvo (delta_hat from eq. (2))
// delta = [0.75, 0.50, 0.25], beta = 0.99
// dhat_i = delta_i*(1-beta*(1-delta_i))/(1-beta*delta_i*(1-delta_i))
DHAT1   = 0.693;
DHAT2   = 0.336;
DHAT3   = 0.079;

// Labour shares (alpha_i = 1 - sum_j(omega_H_ij + omega_F_ij))
ALPHA1  = 0.70;   // Resource:  1 - 0.00 - 0.30
ALPHA2  = 0.70;   // Manuf.:    1 - 0.20 - 0.10
ALPHA3  = 0.70;   // Services:  1 - 0.25 - 0.05

// Domestic I-O shares
OH21    = 0.20;   // Manuf. <- Resource (domestic)
OH32    = 0.25;   // Services <- Manuf. (domestic)

// Total import shares per sector
OF1     = 0.30;   // Resource
OF2     = 0.10;   // Manufacturing
OF3     = 0.05;   // Services

// Final consumption shares
BH1     = 0.05;
BH2     = 0.15;
BH3     = 0.80;
BF_TOT  = 0.10;

// DC index weights  w_i propto lambda_D_i*(1-dhat_i)/dhat_i
// lambda_D = [0.12, 0.35, 0.80] (domestic supplier centrality)
// raw weights:
//   w1 = 0.12*(1-0.693)/0.693 = 0.0532
//   w2 = 0.35*(1-0.336)/0.336 = 0.6916
//   w3 = 0.80*(1-0.079)/0.079 = 9.328
//   total = 10.073
WDC1    = 0.0053;
WDC2    = 0.0686;
WDC3    = 0.9261;

// Open economy
PSI      = 0.020;
THETA_S  = 2.00;
KAP_EX   = 0.85;
KAP_IM   = 0.45;
NU_D     = 0.30;

// Policy
PHI_PI   = 1.50;
PHI_Y    = 0.50;
PHI_S    = 0.30;   // managed float only

// Shock persistence
RHO_A    = 0.90;
RHO_PF   = 0.85;
RHO_D    = 0.80;

// ---------------------------------------------------------
// MODEL BLOCK
// ---------------------------------------------------------
model(linear);

//----------------------------------------------------------
// [1] Sectoral Phillips curves  (eqs. 3-5 in paper)
// pi_it = beta*(1-dhat_i)*E[pi_{i,t+1}]
//       + dhat_i*[alpha_i*w + omega_H_ij*p_j - p_i(-1)
//                + omega_F_i*(s+pF) - a_i]
// Note: p_jt = p_j(-1) + pi_j  is substituted inline
//----------------------------------------------------------

// Resource (no domestic inputs)
pi1 = BETA*(1-DHAT1)*pi1(+1)
    + DHAT1*(   ALPHA1*w_hat
              + OF1*(s_hat + pF)
              - a1
              - p1(-1)
            );

// Manufacturing (buys from Resource domestic)
pi2 = BETA*(1-DHAT2)*pi2(+1)
    + DHAT2*(   ALPHA2*w_hat
              + OH21*(p1(-1) + pi1)
              + OF2*(s_hat + pF)
              - a2
              - p2(-1)
            );

// Services (buys from Manufacturing domestic)
pi3 = BETA*(1-DHAT3)*pi3(+1)
    + DHAT3*(   ALPHA3*w_hat
              + OH32*(p2(-1) + pi2)
              + OF3*(s_hat + pF)
              - a3
              - p3(-1)
            );

//----------------------------------------------------------
// [2] Price level laws of motion  (eq. 6)
//----------------------------------------------------------
p1 = p1(-1) + pi1;
p2 = p2(-1) + pi2;
p3 = p3(-1) + pi3;

//----------------------------------------------------------
// [3] Consumer price level  (eq. 9)
//----------------------------------------------------------
pC_hat = BH1*p1 + BH2*p2 + BH3*p3 + BF_TOT*(s_hat + pF);

//----------------------------------------------------------
// [4] CPI inflation  (eq. 8)
//----------------------------------------------------------
piCPI = BH1*pi1 + BH2*pi2 + BH3*pi3
      + BF_TOT*((s_hat + pF) - (s_hat(-1) + pF(-1)));

//----------------------------------------------------------
// [5] Real wage  (eq. 7, labour supply)
// w_hat = pC_hat + (gamma+varphi)*y_gap
//----------------------------------------------------------
w_hat = pC_hat + (GAMMA + VARPHI)*y_gap;

//----------------------------------------------------------
// [6] DC inflation index  (eq. 10, Proposition 1)
//----------------------------------------------------------
piDC = WDC1*pi1 + WDC2*pi2 + WDC3*pi3;

//----------------------------------------------------------
// [7] IS curve  (eq. 11)
// y_gap = E[y_gap(+1)] - (1/gamma)*(i_hat - E[piCPI(+1)])
//       + nu_D * D_star
// (natural rate r_nat = 0 in deviations from SS)
//----------------------------------------------------------
y_gap = y_gap(+1) - (1/GAMMA)*(i_hat - piCPI(+1)) + NU_D*D_star;

//----------------------------------------------------------
// [8] UIP  (eq. 12)
// i_hat = -psi*bstar + E[s(+1)] - s
//----------------------------------------------------------
i_hat = -PSI*bstar + s_hat(+1) - s_hat;

//----------------------------------------------------------
// [9] NFA accumulation  (eq. 13)
// bstar = (1/beta - psi)*bstar(-1)
//       + kap_EX*(-theta_s*s + D_star)
//       - kap_IM*(pF + y_gap)
// Eigenvalue = 1/beta - psi = 0.990 < 1 (stationary)
//----------------------------------------------------------
bstar = (1/BETA - PSI)*bstar(-1)
      + KAP_EX*(-THETA_S*s_hat + D_star)
      - KAP_IM*(pF + y_gap);

//----------------------------------------------------------
// [10] Policy rule (regime-dependent)
//----------------------------------------------------------
@#if REGIME == "float"
    // Free float: Taylor rule on DC index
    i_hat = PHI_PI*piDC + PHI_Y*y_gap;

@#elseif REGIME == "peg"
    // Hard peg: exchange rate fixed
    s_hat = 0;
    // i_hat determined residually by UIP (already in eq. 8)

@#elseif REGIME == "managed"
    // Managed float: Taylor rule + FX stabilisation
    i_hat = PHI_PI*piDC + PHI_Y*y_gap + PHI_S*s_hat;

@#endif

//----------------------------------------------------------
// [11] Shock processes  (eqs. 15-17)
//----------------------------------------------------------
a1 = RHO_A*a1(-1)  + eps_a1;
a2 = RHO_A*a2(-1)  + eps_a2;
a3 = RHO_A*a3(-1)  + eps_a3;
pF = RHO_PF*pF(-1) + eps_pF;
D_star = RHO_D*D_star(-1) + eps_D;

end;

// ---------------------------------------------------------
// STEADY STATE
// (All log-deviations are zero by construction)
// ---------------------------------------------------------
initval;
    pi1 = 0; pi2 = 0; pi3 = 0;
    piDC = 0; piCPI = 0;
    y_gap = 0; i_hat = 0; s_hat = 0;
    p1 = 0; p2 = 0; p3 = 0;
    pC_hat = 0; w_hat = 0; bstar = 0;
    a1 = 0; a2 = 0; a3 = 0;
    pF = 0; D_star = 0;
end;

// ---------------------------------------------------------
// SHOCKS (unit standard deviations)
// ---------------------------------------------------------
shocks;
    var eps_a1  = 0.01^2;   // 1% TFP shock, sector 1 (Resource)
    var eps_a2  = 0.01^2;   // 1% TFP shock, sector 2 (Manuf.)
    var eps_a3  = 0.01^2;   // 1% TFP shock, sector 3 (Services)
    var eps_pF  = 0.01^2;   // 1% import price shock
    var eps_D   = 0.01^2;   // 1% foreign demand shock
end;

// ---------------------------------------------------------
// SOLUTION AND OUTPUT
// ---------------------------------------------------------
stoch_simul(
    order       = 1,       // linear model
    periods     = 0,       // no simulation; compute IRFs only
    irf         = 40,      // 40-quarter IRF horizon
    graph_format = pdf,
    nograph     = 0,
    noprint     = 0
);

// Print key model moments
disp('=== DC index weights ===');
disp([WDC1 WDC2 WDC3]);

disp('=== Import centrality (network exposure to FX shock) ===');
// M_i = [((I-OmegaH)^{-1}) * OF] -- computed externally from network matrices
// Values: [0.30, 0.16, 0.09] for Resource, Manuf., Services
\end{lstlisting}

\paragraph{(Historical) switching regimes, peg residual, welfare snippet.}
The three paragraphs and Matlab snippet below describe the superseded
linear draft (variable names \texttt{s\_hat}, \texttt{i\_hat}, \texttt{pi1-3}).
In the current nonlinear file the same ideas hold with the new variable
names (\texttt{S}, \texttt{I}, \texttt{PI1-3}): switch regimes via the
\texttt{@\#define REGIME} macro (or let
\texttt{code/run\_all\_regimes.m} do it for all three); under the peg,
\texttt{I} is residual, pinned down by the UIP equation; and the
welfare computation is now \texttt{compute\_welfare()} in
\texttt{code/analysis.py}, reading \texttt{results/variances.csv} and
\texttt{results/network\_objects.csv} exported by
\texttt{run\_all\_regimes.m} -- same formula
(eq.~\ref{eq:welfare} below), but evaluated on the true nonlinear
model's order-2 unconditional variances rather than a linearized
approximation.
\begin{equation}
\mathcal{W}=\tfrac{1}{2}\!\left[(\gamma+\varphi)\,\text{Var}(\tl y_t)
+\sum_i\lambda_{D,i}\varepsilon_i\frac{1-\wh\delta_i}{\wh\delta_i}\,\text{Var}(\pi_{it})\right].
\label{eq:welfare}
\end{equation}

% ════════════════════════════════════════════════════════════════════════════
\section*{References}
% ════════════════════════════════════════════════════════════════════════════

\begin{itemize}\setlength\itemsep{2pt}
\item[\small\textbullet] Gal\'i, J.\ and Monacelli, T.\ (2005). Monetary policy and exchange rate
  volatility in a small open economy. \textit{RES} 72(3), 707--734. 
\item[\small\textbullet] Pasten, E., Schoenle, R.\ and Weber, M.\ (2017). Price rigidity and the
  origins of aggregate fluctuations. NBER WP 23750. 
\item[\small\textbullet] Qiu, Z., Wang, Y., Xu, L.\ and Zanetti, F.\ (2025/26). Monetary policy in
  open economies with production networks. \textit{JME}, forthcoming. CEPR DP19828.
\item[\small\textbullet] Rubbo, E.\ (2024). Networks, Phillips curves, and monetary policy.
  \textit{Econometrica} 91(4), 1417--1455. 
\item[\small\textbullet] Schmitt-Groh\'e, S.\ and Uribe, M.\ (2003). Closing small open economy
  models. \textit{JIE} 61(1), 163--185. 
\item[\small\textbullet] Smets, F.\ and Wouters, R.\ (2003). An estimated dynamic stochastic
  general equilibrium model of the euro area. \textit{JMCB} 35(5), 1123--1162.
  
\end{itemize}

\end{document}
